\section{Introduction}
Visual world models predict the consequences of actions from images. A common
approach rolls a predictor forward, feeding each predicted visual state into
the next step. This makes a long forecast depend increasingly on its own
estimates. The last real observation, however, remains a measured description
of the scene. We ask how a compact model can keep that evidence available
while predicting action-dependent changes.

We propose \textbf{ShiftWM}, an observation-anchored spatial predictor built
from released LeWorldModel components \citep{maes2026lewm}. Its decoder keeps
the last observed feature grid fixed. For each supplied action prefix, it
predicts how to mix features from that grid and adds a learned correction.
The original variant bounds this correction; a controlled variant removes the bound.
The model can therefore change the predicted state without replacing its
observed visual reference with an imagined one.
Figure~\ref{fig:editorial-teaser} contrasts these two routes to a forecast.
Its input gallery also locates the recorded-video transfer tasks and simulation
studies documented in the attached appendix.

\input{generated/editorial/anchoring_teaser_figure.tex}

We evaluate this design through a matched study on recorded DROID robot video.
The proposed model reduces native ten-step endpoint feature error by
\textbf{5.30\%} relative to autoregression and \textbf{3.55\%} relative to an
additive observed-state anchor. The component experiments support the
\emph{action-conditioned mixing package}; they do not establish an incremental
native-error advantage for bounding or the separate support-context path.
We report the complete model and controlled removals under one matched
evaluation protocol.
Separately trained single-image adapters test the same decoder on recorded
IWS pushing and bimanual manipulation tasks, under a distinct forecasting
and selection protocol. A complete reserved comparison of all 36 fixed
predictors finds that the no-$\tanh$ component improves on matched
autoregression across all three tasks.

The contribution is a specific decoder: action-prefix queries mix a fixed
observed feature grid, with a gated identity path and learned correction.
The matched component analysis tests this parameterization, and the real-video
and simulator evaluations delimit its predictive benefits and control failures.
Observation reuse itself has precedents (Section~\ref{sec:editorial-method});
compactness and the temporal backbone are inherited from LeWM. DROID findings are exploratory development results; the
IWS study additionally reports a separate reserved population and its numerical
execution revision. Earlier context-model
studies, including mixed planning and extrapolation outcomes, remain explicitly
separate in the attached appendix and do not establish performance of the
spatial model.

\section{Task and available information}
\label{sec:editorial-task}
Given an observed history and supplied actions, predict future visual features.
The DROID history contains three real images and two past action blocks; the
query contains ten future blocks. Each block groups five native recorded
commands. A frozen encoder maps images into a spatial grid. Withheld future
images supply evaluator targets and never enter support or context inference.
The primary DROID task measures feature forecasting; secondary IWS RGB decoding
and simulator control are separate evaluations, with no physical robot trials.
Figure~\ref{fig:editorial-spatial-method} shows the
model-visible computation; Figure~\ref{fig:editorial-qualitative} shows recorded
observations and measured forecast errors.
